% translation-source: docs/user/latex/chapters/04-raw-profils.tex; sha256: a0c0905dd992d79083ed2f7d606aa1c18b4a5e64831463414261424d5e8d6b39
\chapter{RAW development and profiles}
\section{Choose the starting rendition}
The \emph{RAW Development / Profile} group selects a DCP profile compatible
with the photo. Choices depend on the camera, installed profiles and build.
Check the description of the profile actually used, not just the requested
name in a list.

A camera profile defines a starting rendition: it may change colors, contrast
and apparent brightness before edits. Two applications need not produce
identical pixels from similarly named profiles. This group does not treat
a JPEG as a RAW and apply a camera DCP to it.

\section{Change profiles}
Select a list entry, wait for processing, then compare. The engine can avoid
a complete re-decode when a compatible buffer is available; the interface
indicates this capability. It does not guarantee that every development
change is instantaneous. Commands that could change development are disabled
during some operations.

\section{Automatic white balance and calibration}
The \emph{Automatic white balance} checkbox appears here once. Manual
temperature and tint remain under \emph{Color}. Do not confuse automatic
balance, shooting temperature and the DCP profile.

RAW calibration is an advanced preference, accessed using
\touche{Ctrl + comma}, \emph{Advanced · RAW} tab. Compatible cameras and modes
have default calibration. Check its applied status in the description:
an enabled checkbox does not guarantee that an unknown file has an available
characterization.

\attention{The advanced choice concerns the open image and its sidecar, not
a new global preference for every subsequent photo. Change the calibration
manifest only under a controlled protocol.}

\paragraph{Example.}
To compare profiles, keep edits identical, disable corrections unnecessary
for the test, change only the profile, then inspect skin tones, vivid colors
and highlights. Record white balance for reproducibility.
